% 35

\eocesol{(a)~Tabel dua arah:
\begin{center}\scriptsize
\begin{tabular}{l l c c c}
& \multicolumn{2}{c}{\textit{Berhenti merokok}} &       \\
\cline{2-3}
\textit{Perlakuan}      & Ya       & Tidak        & Total \\
\hline
Koyo + kelompok dukungan   & 40            & 110   & 150   \\
Hanya koyo          & 30            & 120   & 150   \\
\cline{1-4}
Total               & 70            & 230   & 300 \\
\cline{1-4}
\end{tabular}
\end{center}
(b-i)~$E_{baris_1, kolom_1} = \frac{(total~baris~1)\times(total~kolom~1)}{total~tabel} = 35$.
Nilai ini lebih rendah daripada nilai teramati. \\
(b-ii)~$E_{baris_2, kolom_2} = \frac{(total~baris~2)\times(total~kolom~2)}{total~tabel} = 115$.
Nilai ini lebih rendah daripada nilai teramati.}

\end{multicols}
\newpage
\begin{multicols}{2}

% 37

\eocesol{$H_0$: Pendapat lulusan dan bukan lulusan perguruan tinggi tidak berbeda mengenai 
pengeboran minyak dan gas alam di lepas pantai 
California. $H_A$: Pendapat mengenai pengeboran minyak dan gas alam 
di lepas pantai California berasosiasi dengan kepemilikan 
gelar perguruan tinggi.
\begin{align*}
&E_{baris~1, kolom~1} = 151.5 && E_{baris~1, kolom~2} = 134.5 \\
&E_{baris~2, kolom~1} = 162.1 && E_{baris~2, kolom~2} = 143.9 \\
&E_{baris~3, kolom~1} = 124.5 && E_{baris~3, kolom~2} = 110.5
\end{align*}
Saling bebas: Kedua sampel bersifat acak, tidak berkaitan, dan masing-masing kurang 
dari 10\% populasi, sehingga asumsi saling bebas antar-pengamatan 
masuk akal. Ukuran sampel: Semua cacah harapan sekurang-kurangnya 5.
$\chi^2 = 11.47$, $df = 2$ $\to$ nilai-p = 0.003.
Karena nilai-p $< \alpha$, kita menolak $H_0$. Terdapat bukti kuat 
bahwa dukungan terhadap pengeboran lepas pantai berasosiasi dengan 
kepemilikan gelar perguruan tinggi.}
